% page 41 du fac-simile (image p-040 du scan)
% bloc 175.3 x 254.9 mm ; marge gauche 26.7 mm
\pgc{5.080mm}{0.000mm}{
\l{}
\l{}
\l{}
\l{}
\l{}
\l{\sou{aparta}. Quaze separita, izolita, sola, quankam permananta kom parto}
\l{\cel{3}di ensemblo. - DEFIS.}
\l{}
\l{\sou{apartamento}. Parto de domo, qua konsistas ye ula quanto de chambri}
\l{\cel{3}e qua esas habiteyo komfortoza. - DEFIRS.}
\l{}
\l{\sou{apartena}r (\sou{netrans}, ad) Esar la proprietajo di ulu. (\sou{Anke metaf.}).}
\l{\cel{3}- EFIS.}
\l{}
\l{\sou{apatio}. Stando di anmo o di mento o\cel{2}di psiko quin nulo povas}
\l{\cel{3}emocigar.-- DEFIRS.}
\l{}
\l{\sou{apelar}.\cel{2}(\sou{netrans., ad ulu, pri ulo})\cel{2}Rekursar a la judiciistaro kom-}
\l{\cel{3}petenta di la grado superiora, por igar reformar verdikto qua}
\l{\cel{3}enuncesis da tribunalo inferiora. - DEFIRS.}
\l{}
\l{\sou{apendico}.\cel{2}Parto qua servas kom parto di bloko precipua. - DEFIS.}
\l{}
\l{\sou{apendicito}. (\sou{patol.}) Inflamuro di la apendico cekala vermoforma.}
\l{\cel{3}- DEFIS.}
\l{}
\l{\sou{apene}.\cel{2}Preske ne- . - FIS.}
\l{}
\l{\sou{aperceptar}. (\sou{trans}.)(\sou{filoz}.)Intuico, fakultato od ago quakonsistas}
\l{\cel{3}ye intelektar quik, per koncio, ta o ca ideo o verajo. (En la}
\l{\cel{3}filozofio da Leibniz, ta vorto signifikis la percepto juntata}
\l{\cel{3}a reflekto).}
\l{}
\l{\sou{apertar}. (\sou{trans, por}). Igar acesebla per deprenar to quo\cel{2}kluzigas.}
\l{\cel{3}- DEFIS.}
\l{}
\l{\sou{apetitar}. (\sou{trans}.)Tendencar a to quo satisfacos sua bezoni. (\sou{Anke}}
\l{\cel{3}\sou{metaf}.). - DEFIRS.}
\l{}
\l{\sou{apikala}. Dicesas pri la parto qua esas la somito di organo.}
\l{}
\l{\sou{apie,}\cel{1}\sou{(botaniko)}\cel{1}Planto umbelifera, di qua un speco divenis, per}
\l{\cel{3}kultivado, celerio. - L.\cel{1}\sou{apium}.\cel{2}DEFIRS.}
\l{}
\l{\sou{api-pomo}. (\sou{bot.}) Mikra pomo krakanta, di qua un latero esas blank-}
\l{\cel{3}atra, e la altra, reda. - L.\cel{1}\sou{malum appianum}.}
\l{}
\l{\sou{aplanatika}. (\sou{Opt.})\cel{2}En qua ne eventas aberaco sferesala. - DEF.}
\l{}
\l{\sou{aplanata}. (\sou{fotogr}.) Tipo di objektali rektalinea e simetra.- DEF.}
\l{}
\l{\sou{aplastar}. (\sou{trans}.)I. Platigar e deformar (korpo) per sinkar ula}
\l{\cel{3}parto per shoko violentoza o per granda preso. - II. (\sou{Metaf}.)}
\l{\cel{3}Igar sukombar ad pezozajo tro grava, a kargajo tro grava ;}
\l{\cel{3}igar sukombar en lukto per efikigar forco a qua ne esas posibla}
\l{\cel{3}rezistar. - FIS.}
\l{}
\l{\sou{aplaudar}. (\sou{trans}.)Intershokigar la parto interna di la manui kom}
\l{\cel{3}manifesto di sua granda admiro. - DEFIRS.}
}
